%% -----------------------------------------------------------
\section{Policy Implications}
%% -----------------------------------------------------------

The evidence supports five urgent policy interventions.

\textit{Mandatory pre-deployment safety standards} should require standardized testing for crisis recognition, cultural bias, neurotypical bias, accessibility, data privacy, and adverse event monitoring before any AI tool is deployed for vulnerable populations. The current absence of such requirements, documented by Smith et al.~\cite{smith_2023} and reinforced by the \gls{neda} case~\cite{neda_2023}, represents a structural failure of regulatory oversight.

\textit{Community co-design certification} should mandate documented community engagement processes, cultural adaptation evidence, accessibility testing with disabled users, neurodivergent user testing, and ongoing advisory boards for any AI tool claiming to serve marginalized populations.

\textit{Transparency and informed consent} should provide users with clear information about what chatbots address, what evidence supports effectiveness, what limitations apply, when human support is necessary, how data is used, and what crisis escalation exists. Palmer et al.~\cite{palmer_2025} propose a voluntary certification program as one practical mechanism for standardizing these disclosures.

\textit{Hybrid model prioritization} for high-risk populations should require human oversight for any AI deployment serving populations with elevated suicide risk or acute mental health vulnerability, with particular attention to neurodivergent and \gls{lgbtq} youth.

\textit{Accessibility as infrastructure} should require platforms to meet \gls{wcag} standards and demonstrate usability across assistive technologies before deployment, treating accessibility as a prerequisite for service provision~\cite{khan_seo_2025}.
